% !TEX root = main.tex
\subsection{Transforming PLONK's Constraints to Quadratic Constraints}
\label{subsection:transforming-plonk-arithmetization-to-quadratic-constraints}

\noindent\textbf{PLONK's constraint system.}
PLONK's arithmetization~\cite{EPRINT:GabWilCio19} encodes $\Circuit$
via selectors $\{s_i^\sfleft, s_i^\sfright, s_i^\sfmult, s_i^\sfconst\}_{i = 1}^{\ngate}$
and a permutation $\varphi \colon [\nw] \to [\nw]$ encoding wire-sharing.
Let $\ellplonk \coloneqq 4\ngate + \nw$.
Following the formulation of Tang, Pham, and Ngo~\cite{EPRINT:TanPhaNgo24},
these are collected into the \emph{PLONK description}:
\begin{equation}\label{equation:plonk-circuit-description}
\plonkdescription =
\bigl(
(s_i^\sfleft)_{i = 1}^{\ngate}
\;\|\;
(s_i^\sfright)_{i = 1}^{\ngate}
\;\|\;
(s_i^\sfmult)_{i = 1}^{\ngate}
\;\|\;
(s_i^\sfconst)_{i = 1}^{\ngate}
\;\|\;
(\varphi(i))_{i = 1}^{\nw}
\bigr)
\;\in\; \FF^{\ellplonk},
\end{equation}
This vector fully encodes the circuit's gate structure and wiring.

For any input $\bx \in \FF^\nin$, evaluating $\Circuit$ yields gate-value sequences
$\ba = (a_i)_{i = 1}^{\ngate}$,
$\bb = (b_i)_{i = 1}^{\ngate}$,
$\bc = (c_i)_{i = 1}^{\ngate}$,
and output $\bcout = (c_i)_{i = \ngate - \nout + 1}^{\ngate}$ for some $\nout \leq \ngate$.
The \emph{extended witness} is:
\begin{equation}\label{equation:plonk-tilde-w}
\tildebw = (\bx \;\|\; \ba \;\|\; \bb \;\|\; \bc)
= (w_i)_{i = 1}^{\nw} \in \FF^{\nw}.
\end{equation}
Together with $\plonkdescription$, $\tildebw$ carries all information
needed to verify $\Circuit$'s evaluation.

For random challenges $\alpha, \beta \in \FF$, the output $\bcout$ is correct
if and only if:
\begin{equation}\label{equation:plonk-system-of-constraints-with-challenges}
\begin{cases}
\begin{aligned}
	&s_i^\sfleft \cdot w_{\nin + i} + s_i^\sfright \cdot w_{\nin + \ngate + i}
	\\
	&\quad + s_i^\sfmult \cdot (w_{\nin + i} \cdot w_{\nin + \ngate + i})
	\\
	&\quad + s_i^\sfconst - w_{\nin + 2 \ngate + i} = 0
\end{aligned}
& \forall i \in [\ngate],
\\[15pt]
\displaystyle
\prod_{i = 1}^{\nw}(\alpha + \beta \cdot i + w_i)
=
\prod_{i = 1}^{\nw}(\alpha + \beta \cdot \varphi(i) + w_i),
\quad \alpha, \beta \uniformly \FF.
\end{cases}
\end{equation}
The first equation encodes gate constraints;
the second encodes wire consistency via a grand product check.
This system can be viewed as a circuit $\UniversalCircuitplonk$
taking $(\plonkdescription, \tildebw, \alpha, \beta)$ as inputs
and verifying whether $\Circuit(\bx) = \bcout$.
The soundness error is at most $\nw / |\FF|$.
\ifcameraready
\else
See Appendix~\ref{subappendix:preliminaries-plonk-arithmetization}
for further background.
\fi

\medskip
\noindent\textbf{Reduction to quadratic constraints.}
The gate constraint in~\eqref{equation:plonk-system-of-constraints-with-challenges}
contains the cubic term $s_i^\sfmult \cdot (w_{\nin + i} \cdot w_{\nin + \ngate + i})$.
To reduce to purely quadratic constraints, introduce auxiliary variables:
$\hatw_i \coloneqq w_{\nin + i} \cdot w_{\nin + \ngate + i}$ for each $i \in [\ngate]$.
For the grand product, define running products:
\begin{align*}
\tildew_1 &\coloneqq \alpha + \beta + w_1,
& w'_1 &\coloneqq \alpha + \beta \cdot \varphi(1) + w_1,
\\
\tildew_i &\coloneqq \tildew_{i - 1} \cdot (\alpha + \beta \cdot i + w_i)
& w'_i &\coloneqq w'_{i - 1} \cdot (\alpha + \beta \cdot \varphi(i) + w_i)
\end{align*}
for all $i \in [2, \nw]$.
The grand product holds if and only if $\tildew_{\nw} = w'_{\nw}$.
System~\eqref{equation:plonk-system-of-constraints-with-challenges}
is therefore equivalent to the following purely quadratic system:
\begin{equation}\label{equation:plonk-quadratic-constraints-with-challenges}
\begin{cases}
\hatw_i - w_{\nin + i} \cdot w_{\nin + \ngate + i} = 0
& \!\!\forall i \in [\ngate],
\\
s_i^\sfleft \cdot w_{\nin + i}
+ s_i^\sfright \cdot w_{\nin + \ngate + i}
+ s_i^\sfmult \cdot \hatw_i
+ s_i^\sfconst - w_{\nin + 2 \ngate + i} = 0
& \!\!\forall i \in [\ngate],
\\
\alpha + \beta + w_1 - \tildew_1 = 0,
%\quad
\alpha + \beta \cdot \varphi(1) + w_1 - w'_1 = 0,
%\quad
\tildew_{\nw} = w'_{\nw},
\\
\tildew_{i - 1} \cdot (\alpha + \beta \cdot i + w_i) - \tildew_i = 0
& \!\!\forall i \in [2, \nw],
\\
w'_{i - 1} \cdot (\alpha + \beta \cdot \varphi(i) + w_i) - w'_i = 0
& \!\!\forall i \in [2, \nw].
\end{cases}
\end{equation}
%Crucially, s
System~\eqref{equation:plonk-quadratic-constraints-with-challenges}
has the same structure for any arithmetic circuit of the same gate count,
independent of the specific circuit.
This uniformity is what enables hidden-circuit satisfiability proofs:
the constraints do not expose the circuit's topology.

All witnesses supporting verification of~\eqref{equation:plonk-quadratic-constraints-with-challenges}
are collected in:
\begin{equation}\label{equation:plonk-ww}
\wit = \bigl(
\plonkdescription \;\|\;
\bx \;\|\; \ba \;\|\; \bb \;\|\; \bc \;\|\;
(\hatw_i)_{i = 1}^{\ngate} \;\|\;
(\tildew_i)_{i = 1}^{\nw} \;\|\;
(w'_i)_{i = 1}^{\nw}
\bigr).
\end{equation}
This vector is fully determined by $\alpha$, $\beta$, and $(w_i)_{i = 1}^{\nw}$.

\medskip
\noindent\textbf{Two-phase witness structure.}
The vector $\wit$ splits into two parts:
$\wit = (\witpre \;\|\; \witsuf)$, where
\begin{align*}
\witpre &= \bigl(\plonkdescription \;\|\; \bx \;\|\; \ba \;\|\; \bb \;\|\; \bc \;\|\;
(\hatw_i)_{i = 1}^{\ngate}\bigr)
\quad\text{is determined without knowing } \alpha, \beta;
\\
\witsuf &= \bigl((\tildew_i)_{i = 1}^{\nw} \;\|\; (w'_i)_{i = 1}^{\nw}\bigr)
\quad\text{is determined only after } \alpha, \beta \text{ are known.}
\end{align*}
This split is critical for the ZKBoo-based instantiation of Section~\ref{section:overview-lattice-sms}:
the prover first commits to $\witpre$,
derives $\alpha, \beta$ from a random oracle,
then computes and commits to $\witsuf$.

In summary, any hidden arithmetic circuit $F$ can now be verified
via a fixed public system of quadratic constraints of size
$\BigO(\nin + \nadd + \nmult)$,
with the circuit's structure encoded entirely in the witness $\wit$.


\medskip
\noindent\textbf{Sizes.}
Since $\nw = \BigO(\nin + \ngate) = \BigO(\nin + \nadd + \nmult)$,
the witness vector $\wit$ has $\BigO(\nin + \nadd + \nmult)$ field elements,
and system~\eqref{equation:plonk-quadratic-constraints-with-challenges}
has $\BigO(\nin + \nadd + \nmult)$ multiplications and additions over $\FF$.
